% 附录

% \zufeAppendix 会自动把章节编号切换为字母，并复位图、表、公式计数器，
% 因此附录中的编号为：图 A1、表 A1、公式 (A1)，与正文分开另行编号。
% 多个附录依次调用即可：\zufeAppendix{附录内容的名称} → 附录A、附录B……
% 不写附录时，把 chapters/mainbody.tex 中引入本文件的 \input 注释掉。
\zufeAppendix{附录内容的名称}

% 附录相关内容…

% % 这里示范一下添加多个附录的方法：

% \section{\LaTeX 环境的安装}
% \LaTeX 环境的安装。

% \section{使用说明}
% 使用说明。

\textcolor{blue}{以下内容可放在附录之内：
1．正文内过于冗长的公式推导；
2．方便他人阅读所需的辅助性数学工具或表格；
3．重复性数据和图表；
4．论文使用的主要符号的意义和单位；
5．程序说明和程序全文;
6．调研报告。
}

\textcolor{blue}{这部分内容可省略（如果省略，删去此页）。阅后删除此段。}
